% Apply-ready replacement fragments for the supplied 20-page manuscript.
% This is NOT a replacement main manuscript: copy each block at its anchor.
% Exact anchors, label placement, validation and bibliography mapping are in
% theory-edits.md. Requires amsmath/amssymb; preserve existing theorem labels.

% M01: replace §2.1 "where D(X) is the intervention domain."
where $D(X)$ is the intervention domain. We use fixed nonempty per-variable
intervention domains $D_v$ and set $D(X)=\prod_{v\in X}D_v$.
Thus the restriction of an admissible assignment to any subset of its
intervened variables remains admissible. These domains need not contain
all naturally attainable values; the stronger completion-support condition
is stated separately in Definition~1.

% M02: replace §4.2 sentence beginning "The MIS-reduction result ..."
Appendix~A shows that MIS reduction preserves the population optimum when
the null intervention is included. Our action families and implemented
arm-based optimizers exclude that intervention. Equality between $V(\Pi)$
and the best retained nonempty quotient-MIS value therefore needs an
additional condition, such as a retained action with population value at
most $\mathbb E_M[Y]$. This conclusion requires a graph compatible with the
interventional distributions being evaluated; it does not apply to an
arbitrary misspecified supplied graph. Proposition~2 concerns $\mathcal U(\Pi)$
directly and does not require MIS reduction.

% M02: replace all of Lemma 1, its proof and the "Consequently ..." paragraph
% in Appendix A.3; preserve the following X1 -> X2 -> Y example.
\paragraph{Lemma 1 (MIS reduction, with the null intervention).}
Let $G$ be the latent-projection ADMG of the acyclic SCM $M$ whose
interventional distributions are evaluated, with manipulable variables
$\mathcal X$ and target $Y$. Assume the product-domain convention of
Section~2.1. For $X\subseteq\mathcal X$, put
\[
 X'=X\cap\operatorname{An}_{G_{\overline X}}(Y).
\]
Then (i) for every $x\in D(X)$,
\[
 P_M(Y\mid\operatorname{do}(X=x))
 =P_M(Y\mid\operatorname{do}(X'=x|_{X'}));
\]
(ii) if $X'\ne\varnothing$, then $X'\in\operatorname{MIS}(G)$;
and (iii), defining $D(\varnothing)=\{\varnothing\}$ and
$f_{\varnothing}(\varnothing)=\mu_{\varnothing}:=\mathbb E_M[Y]$,
\[
 \min_{X\in\{\varnothing\}\cup\operatorname{MIS}(G)}\mu^\star(X)
 =\min_{X\subseteq\mathcal X}\mu^\star(X).
\]
Here MIS retains the manuscript's nonempty-set convention, and the
attainment assumptions of Section~4.2 apply.

\paragraph{Proof.}
Write $R=X\setminus X'$. If some $r\in R$ had a directed path to $Y$
in $G_{\overline {X'}}$, that path could not enter $X'$. Let $r^\star$
be the last member of $R$ on the path. The remaining path from $r^\star$
to $Y$ enters no member of $X$, so it also exists in $G_{\overline X}$,
contradicting the definition of $R$. Hence restoring the equations for
$R$ while holding $X'$ fixed cannot affect $Y$. Coupling the two
interventions with the same exogenous disturbances proves (i), including
when $X'$ is empty. Domain projection makes $x|_{X'}$ admissible.
For each $v\in X'$, its directed path to $Y$ in $G_{\overline X}$ also
exists in $G_{\overline {X'}}$; thus (ii) follows from the MIS criterion.
Every candidate value on the right side of (iii) is therefore represented
on the left, and the reverse inequality follows by action-family inclusion.

For a valid quotient of this true graph, let
$\mathcal E_\Pi=\operatorname{ES}_\Pi(G_\Pi)$,
$W(\Pi)=\min_{X\in\mathcal E_\Pi}\mu^\star(X)$, and set
$W(\Pi)=+\infty$ if $\mathcal E_\Pi$ is empty. All clusters eligible for
intervention are fully manipulable. Apply the same last-removed-ancestor
argument at cluster level: after the retained clusters are fixed, a
fine directed path from any removed cluster to $Y$ would project to a
directed quotient path, a contradiction. Thus removing those whole-cluster
interventions preserves the response, even with latent confounding, and
\[
 \min\{\mu_{\varnothing},W(\Pi)\}
 =\min\{\mu_{\varnothing},V(\Pi)\}.
\]
Since $\mathcal E_\Pi\subseteq\mathcal U(\Pi)$, we also have
$V(\Pi)\le W(\Pi)$. If $W(\Pi)\le\mu_{\varnothing}$, the displayed
equality forces $V(\Pi)=W(\Pi)$: a strictly smaller $V(\Pi)$ would make
the right side strictly smaller than the left. This sufficient condition
is not assumed for arbitrary benchmarks or wrong supplied graphs.

For a counterexample without this condition, take independent standard
normal $U,U_Z$ and set $A=0.1U$, $Z=U_Z$, $Y=A^2$, with manipulable
variables $\{A,Z\}$, $D_A=[1,2]$ and $D_Z=[-1,1]$.
The graph has only the edge $A\to Y$. Its only nonempty MIS is $\{A\}$,
whose optimal value is $1$, whereas intervening on $Z$ alone gives
$\mathbb E[A^2]=0.01$. Hence at singleton resolution
$V=0.01<W=1$. Reducing $\{Z\}$ gives the null intervention and
$\min\{0.01,1\}=\min\{0.01,0.01\}$, as asserted by the corrected lemma.
Observation-only rounds in our implementation update the prior and retain
the existing interventional incumbent; they do not make the null
intervention eligible as a recommendation.

% M03: replace the mechanism display and its following explanatory paragraph
% in §3.2, ending before "Multi-variable clusters are modeled ...".
\[
 K_C(dv_C\mid v_{\operatorname{pa}(C)}),
\]
and composes them along the C-DAG. Here $v_C$ is the vector of member
variables and $\operatorname{pa}(C)$ is the set of parent clusters.
The structural kernel $K_C$ is induced by composing the member-level SCM
mechanisms within $C$ at the specified parent-cluster values. Under causal
sufficiency and a valid partition, it is a version of the observational
conditional, agreeing almost surely in the observational parent distribution.
At other parent values the notation refers to the SCM-defined kernel;
observational data alone need not determine its extension, even at
individual zero-probability values in the support. GP extrapolation beyond
observed support is a modeling assumption. Joint cluster mechanisms preserve
within-cluster dependence.

% M03: replace Proposition 4 statement and proof in Appendix A.1.
\paragraph{Proposition 4 (Soundness of structural cluster kernels).}
Assume an acyclic causally sufficient SCM with mutually independent
node disturbances and a valid partition $\Pi$. For
$\varnothing\ne S\subseteq\Pi_{\mathcal X}$ and $X=\bigcup_{C\in S}C$,
for every $x\in D(X)$,
\[
 P_M(dv_{\mathcal V\setminus X}\mid\operatorname{do}(X=x))
 =\left.\prod_{C\in\Pi\setminus S}
 K_C(dv_C\mid v_{\operatorname{pa}(C)})\right|_{X=x}.
\]
The product is a composition of probability kernels in quotient
topological order; intervened coordinates are fixed constants.

\paragraph{Proof.}
Each node equation and its independent disturbance define a structural
kernel at every parent assignment. For each cluster, compose its member
kernels in their internal topological order, holding parent-cluster values
fixed, to obtain $K_C$. A whole-cluster intervention deletes every member
kernel of each intervened cluster and fixes its coordinates. Grouping the
remaining kernels gives the displayed law. Quotient validity permits an
acyclic cluster order. Since cluster disturbances are independent of their
parent clusters, $K_C$ is a version of the observational conditional.
The structural construction, rather than a choice of observational
conditional on null events, defines its values elsewhere.

% M04: replace §3.1's two singleton sentences (ending "Appendix A.").
At the singleton partition, the graph and intervention construction reduce
to the fine graph and its nonempty MIS. QCBO's prior construction agrees
arm by arm with the matched CBO implementation when it uses the same
estimators, observation policy and numerical settings. This is not an
identity claim for every original CBO MIS/POMIS configuration.
Appendix~A distinguishes construction identity from executable identity.

% M04: replace §4.1's "At the singleton partition, they recover ..." sentence.
At the singleton partition, the quotient graph and action construction
reduce to their fine counterparts. Proposition~3 separates this structural
identity from the matching implementation conditions needed for numerical
baseline equivalence.

% M04: replace Proposition 3 statement and proof in Appendix A.1.
\paragraph{Proposition 3 (Singleton construction identity).}
At $\Pi=\Pi_{\mathrm{fine}}$, quotienting preserves the graph and the QCBO
exploration set is $\operatorname{MIS}(\widehat G)$. With the same
identification and estimation routines, its prior construction agrees
arm by arm with matched nonempty-MIS CBO. The dynamic quotient graphs and
the structural cluster factorization likewise reduce to their fine
counterparts. Equality of instantiated optimizers additionally requires
matching estimators, GP parameterization, acquisition, initialization,
observation policy, numerical settings and random-number streams.

\paragraph{Proof.}
At singleton resolution $\chi$ is one-to-one, no graph edge is changed,
and unions of manipulable singleton clusters are exactly the subsets of
$\mathcal X$. Equation~(6) therefore becomes Equation~(2).
Identification queries coincide, and matched numerical estimators give
matched priors. The same argument applies to dynamic slice and transition
graphs, while a singleton structural kernel is just its node kernel.
These facts prove construction identity. They imply equality of
optimizer trajectories only when the additional implementation inputs
and routines match, in which case induction over updates applies.

Our current QMCBO implementation learns the observation-noise parameter
of a singleton \texttt{SingleTaskGP}; the pinned upstream MCBO uses
\texttt{FixedNoiseGP} with the simulator-supplied additive-noise variance.
Consequently, we do not claim numerical singleton equivalence between
these implementations. This distinction leaves the structural statement
and quotient-preserving graph invariance unchanged.

% M17: replace the named sentence in §5.2; add labels as directed in .md.
Figure~\ref{fig:controlled-misspecification}(a) shows the objective-value
consequence; Table~\ref{tab:controlled-complete} reports cumulative
incumbent regret.
% M17: replace the named sentence in §5.3.
Figure~\ref{fig:controlled-misspecification}(b) shows the slower objective
improvement; Table~\ref{tab:controlled-complete} reports the corresponding
regret increase.

% M21: replace §6 final CAMAB comparison from "A setting related ..."
CAMAB methods relate decision problems through causal abstractions
\cite{zennaro2024camab,dyer2025atucb}. In particular, AT-UCB uses an
abstract model to eliminate candidate actions before learning in the base
model and bounds the relation between optimal rewards across abstraction
levels \cite{dyer2025atucb}. Our setting fixes a supplied quotient and
restricts interventions to whole clusters. Our contribution is optimizer
invariance to quotient-preserving graph changes and the associated
action-restriction and completion analysis. Abstraction-based action
elimination is not claimed as new.

% M21: insert in §6 after the model-based discussion (or its first paragraph).
MCBO already accommodates vector-valued mechanisms
\cite{sussex2023mcbo}. QMCBO replaces fine mechanisms by mechanisms on a
supplied quotient and enforces whole-cluster intervention targets;
vector-output modeling itself is not our novelty.

% M21: insert §4.4 after "These conditions are sufficient but not necessary."
Original CBO also discusses natural mediator responses outside their
intervention domains \cite{aglietti2020cbo}. Our action-family value
$\Delta(\Pi)$ formalizes the restriction imposed by a chosen partition,
while the completion result supplies sufficient conditions for that
restriction to be lossless.
